\SourcePage{485}{488}
\section{Vibrational energy levels}

In a quantum-mechanical treatment, the molecule's vibrational energies are given by the eigenvalues of the Hamiltonian. The Hamiltonian has the form

\begin{equation}
 \widehat H^{(v)}=\frac12\sum_\alpha\sum_{i=1}^{f_\alpha}
 \left(\widehat P_{\alpha i}^{2}+\omega_\alpha^2Q_{\alpha i}^2\right),
 \tag{101.1}
\end{equation}
where $\widehat P_{\alpha i}=-i\hbar\,\partial/\partial Q_{\alpha i}$ are the
momentum operators corresponding to the normal coordinates
$Q_{\alpha i}$. Since this Hamiltonian separates into a sum of independent
terms (the expressions in parentheses), the energy levels are given by the
sums
\begin{equation}
 E^{(v)}=\hbar\sum_\alpha\omega_\alpha\sum_i\left(v_{\alpha i}+\frac12\right)
 =\sum_\alpha\hbar\omega_\alpha\left(v_\alpha+\frac{f_\alpha}{2}\right),
 \tag{101.2}
\end{equation}
where $v_\alpha=\sum_i v_{\alpha i}$ and $f_\alpha$ is the multiplicity of the
frequency $\omega_\alpha$. The wave

\SourcePage{486}{489}
functions, in turn, are products of the corresponding wave functions of
linear harmonic oscillators:
\begin{equation}
 \psi=\prod_\alpha\psi_\alpha,
 \tag{101.3}
\end{equation}
where
\begin{equation}
 \psi_\alpha=\mathrm{const}\cdot\exp\left(-\frac12c_\alpha^2\sum_iQ_{\alpha i}^2\right)
 \prod_iH_{v_{\alpha i}}(c_\alpha Q_{\alpha i});
 \tag{101.4}
\end{equation}
$H_v$ denotes the Hermite polynomial of degree $v$, and
$c_\alpha=\sqrt{\omega_\alpha/\hbar}$.

If some of the frequencies $\omega_\alpha$ are multiple, the vibrational
energy levels are degenerate in general. The energy (101.2) depends only on the
sum $v_\alpha=\sum_i v_{\alpha i}$. The degeneracy of a level is therefore the
number of ways in which a given set of the numbers $v_\alpha$ can be formed
from the numbers $v_{\alpha i}$. For a single number $v_\alpha$, this number of
ways is\footnote{This is the number of ways to distribute $v_\alpha$ balls
among $f_\alpha$ boxes.}
\[
 \frac{(v_\alpha+f_\alpha-1)!}{v_\alpha!(f_\alpha-1)!}.
\]
Thus the total degeneracy is
\begin{equation}
 \prod_\alpha\frac{(v_\alpha+f_\alpha-1)!}{v_\alpha!(f_\alpha-1)!}.
 \tag{101.5}
\end{equation}
For twofold frequencies the factors in this product are $v_\alpha+1$, and for
threefold frequencies they are $(1/2)(v_\alpha+1)(v_\alpha+2)$.

It must be kept in mind that this degeneracy exists only while the vibrations
are treated as purely harmonic. If terms of higher degree in the normal
coordinates are included in the Hamiltonian (anharmonicity of the
vibrations), the degeneracy is lifted in general, though not completely (see
\S~104 for further details).

The wave functions (101.3) belonging to the same degenerate vibrational term
carry a certain representation (reducible in general) of the molecule's
symmetry group. Functions belonging to different frequencies, however,
transform independently of one another. Hence the representation carried by
all the functions (101.3) is the product of the representations carried by
the functions (101.4), so it is enough to consider only the latter.

The exponential factor in (101.4) is invariant under every symmetry
transformation. In the Hermite polynomials, terms of any fixed degree
transform only into one another (a symmetry transformation plainly does not
alter the degree of each term). On the other hand, each Hermite polynomial is
fully determined by its highest-degree term. Consequently, on writing

\SourcePage{487}{490}
\[
 \prod_{i=1}^{f_\alpha}H_{v_{\alpha i}}(c_\alpha Q_{\alpha i})
 =\mathrm{const}\cdot Q_{\alpha1}^{v_{\alpha1}}Q_{\alpha2}^{v_{\alpha2}}\ldots
 Q_{\alpha f_\alpha}^{v_{\alpha f_\alpha}}+\text{terms of lower degree},
\]
it is sufficient to consider the highest-degree term alone.

Functions for which the sum $v_\alpha=\sum_i v_{\alpha i}$ has the same value
belong to the same term. We thus obtain the representation carried by products
of $v_\alpha$ quantities $Q_{\alpha i}$. This is precisely the symmetric
product (see \S~94) of $v_\alpha$ copies of the irreducible representation
carried by the quantities $Q_{\alpha i}$ (L.~Tisza, 1933).

For one-dimensional representations, finding the characters of their
$v$-fold symmetric products is trivial\footnote{Here we use the notation
$\chi_v(G)$ in place of the cumbersome $[\chi^v](G)$.}:
\[
 \chi_v(G)=[\chi(G)]^v.
\]
For two- and three-dimensional representations, the following mathematical
device is convenient\footnote{Applied for this purpose by A.~S.~Kompaneyets
(1940).}. The sum of the squares of the basis functions of an irreducible
representation is invariant under all symmetry transformations. They may
therefore be regarded formally as the components of a two- or
three-dimensional vector, and the symmetry transformations as certain
rotations (or reflections) acting on these vectors. We stress that, in
general, these rotations and reflections have no connection with the actual
symmetry transformations and, for every particular group element $G$, they
also depend on the specific representation under consideration.

Consider two-dimensional representations in more detail. Let $\chi(G)$ be the
character of some group element in a given two-dimensional representation,
with $\chi(G)\ne0$.
Consider the matrix that transforms the components $x$ and $y$ of a two-dimensional vector under a rotation in the plane through an angle $\varphi$.
The sum of its diagonal elements, namely its trace, is $2\cos\varphi$.

By setting
\begin{equation}
 2\cos\varphi=\chi(G),
 \tag{101.6}
\end{equation}
we determine the angle of the rotation that formally corresponds to the
element $G$ in the given irreducible representation. The $v$-fold symmetric
product of the representation has a basis consisting of the $v+1$ quantities
$x^v,x^{v-1}y,\ldots,y^v$. The characters

\SourcePage{488}{491}
of this representation are\footnote{For the calculation it is convenient to
choose the basis functions in the form
\[
 (x+iy)^v,\quad (x+iy)^{v-1}(x-iy),\ldots,(x-iy)^v;
\]
the rotation matrix is then diagonal, and the sum of its diagonal elements
has the form
\[
 e^{iv\varphi}+e^{i(v-2)\varphi}+\ldots+e^{-iv\varphi}.
\]}
\begin{equation}
 \chi_v(G)=\frac{\sin(v+1)\varphi}{\sin\varphi}.
 \tag{101.7}
\end{equation}
The case $\chi(G)=0$ requires separate consideration, since a zero character
corresponds both to a rotation through $\pi/2$ and to a reflection. If
$\chi(G^2)=-2$, the transformation is a rotation through $\pi/2$, and for
$\chi_v(G)$ we obtain
\begin{equation}
 \chi_v(G)=(-1)^{v/2}\frac{1+(-1)^v}{2}.
 \tag{101.8}
\end{equation}
If instead $\chi(G^2)=2$, then $\chi(G)$ must be treated as the character of a
reflection (that is, the transformation $x\to x$, $y\to-y$); in that case
\begin{equation}
 \chi_v(G)=\frac{1+(-1)^v}{2}.
 \tag{101.9}
\end{equation}

Formulas for symmetric products of three-dimensional representations can be
obtained in the same way. The rotation (or reflection) formally corresponding
to a group element in the given representation is readily found with the aid
of Table~7. It is the transformation corresponding to the given $\chi(G)$ in
that one of the isomorphic groups in which the coordinates transform according
to this representation. Thus, for the representation $F_1$ of the groups $O$
and $T_d$, the transformation must be taken from the group $O$, whereas for
the representation $F_2$ it must be taken from the group $T_d$. We shall not
dwell here on the derivation of the corresponding formulas for the characters
$\chi_v(G)$.
